\chapter{除子与可逆层}\label{chap:ag-1}
\section{亚纯函数与 Cartier 除子}\label{sec:ag-1-1}
\begin{numpara}[全分式环]\label{item:ag-1-1-1}
环 $A$ 的元素 $f$ 称为\newterm{ag-regular-element}{正则元素}{regular element}，如果它不是零因子. 记 $S_A$ 为 $A$ 的正则元素全体. 这是 $A$ 的一个``乘性子集''. 令 $Q(A):=S_A^{-1}A$ 为 $A$ 关于 $S_A$ 的局部化，称之为 $A$ 的\newterm{ag-total-ring-of-fractions}{全分式环}{total ring of fractions}.
典范态射 $A\longrightarrow Q(A)$ 总是单射.
例：
\begin{itemize}
\item 如果 $A$ 是整环，则 $S_A=A\setminus\{0\}$，$Q(A)$ 是 $A$ 的分式域.
\item 如果 $A$ 是 Noether 环，则 $Q(A)$ 是 Artin 环，因而是有限多个 Artin 局部环的乘积. 它的极大理想对应于 $A$ 的极小素理想，即 $\operatorname{Spec}(A)$ 各不可约分支的泛点. 如果 $A$ 还既约，则 $Q(A)$ 也既约，因而是若干个域的乘积.
\item 如果 $A$ 是 Artin 环，则 $S_A=A^\times$，且 $Q(A)=A$.
\item $Q(A\times B)=Q(A)\times Q(B)$.
\end{itemize}
\end{numpara}
\begin{numpara}[亚纯函数层]\label{item:ag-1-1-2}
设 $X$ 是概形，$\mathcal O_X$ 是其结构环层. 预层 $U\mapsto S_{\mathcal O_X(U)}$ 实际上是层. 这是因为 $\mathcal O_X(U)$ 的截面 $f$ 正则，当且仅当对每个 $x\in U$，芽 $f_x$ 在 $\mathcal O_{X,x}$ 中正则. （回忆态射 $A\longrightarrow\prod_{\mathfrak p\in\operatorname{Spec}(A)}A_{\mathfrak p}$ 是忠实平坦的，且 $f$ 正则等价于乘以 $f$ 的映射是单射.）
预层 $U\mapsto Q(\mathcal O_X(U))$ 未必是层，记其相伴层为 $\mathcal M_X$，称为 $X$ 上的\newterm{ag-meromorphic-function-sheaf}{亚纯函数层}{sheaf of meromorphic functions}. 这是一个环层，并带有单态射
\begin{align*}
\mathcal O_X\hookrightarrow\mathcal M_X.
\end{align*}
由构造，$\mathcal M_X$ 在 $X$ 的开集 $U$ 上的限制与 $\mathcal M_U$ 相同，且 $\mathcal M_X$ 在点 $x$ 处的茎 $\mathcal M_{X,x}$ 与 $Q(\mathcal O_{X,x})$ 相同.
例：如果 $X$ 是整概形，则 $\mathcal M_X$ 是常值层，其茎是 $X$ 的泛点处的函数域 $\kappa(x)$.
\end{numpara}
\begin{numpara}[亚纯函数及其定义域]\label{item:ag-1-1-3}
记 $\mathcal M(X):=\Gamma(X,\mathcal M_X)$，称为 $X$ 上的亚纯函数环. 设 $f\in\mathcal M(X)$. 存在最大的（非空）开集 $U$，使得 $f|_U\in\mathcal O_X(U)\subset\mathcal M_X(U)$. 称此开集为 $f$ 的\newterm{ag-domain-of-definition}{定义域}{domain of definition}，记作 $\operatorname{dom}(f)$.
例：
\begin{itemize}
\item 如果 $X=\operatorname{Spec}(A)$，其中 $A$ 是 Noether 环，则 $\mathcal M(A)=Q(A)$. 当 $A$ 是整环时，这一点显然，因为 $\mathcal M_X$ 是常值层. 更一般情形的证明可参见 EGA IV (20.2.11) (ii).
\item 如果 $X$ 是整概形，则 $\mathcal M(X)$ 是 $X$ 的泛点处的剩余域（见上文）. 特别地，如果 $X=\mathbb P_k^1$，其中 $k$ 是域，则 $\mathcal M(X)$ 是函数域 $k(T)$，远大于 $Q(\mathcal O_X(X))=k$. 这说明 $U\mapsto Q(\mathcal O_X(U))$ 不是层.
\end{itemize}
\emph{注：函子性.} 对概形态射 $Y\xrightarrow f X$，不能期待 $\mathcal M(X)$ 和 $\mathcal M_X$ 具有一般的函子性（考虑 $Y$ 是 $X$ 的一个点的情形：某些亚纯函数在 $Y$ 上没有定义）. 但是，如果 $f$ 平坦，则归结到仿射情形可构造一个单态射 $f^*\mathcal M_X\longrightarrow\mathcal M_Y$，它延拓自然同构 $f^*\mathcal O_X\xrightarrow\sim\mathcal O_Y$. 事实上，如果环 $B$ 在 $A$ 上平坦，则 $A$ 的每个正则元素在 $B$ 中仍正则，因而有态射
\begin{align*}
Q(A)=S_A^{-1}A\longrightarrow Q(A)\otimes_A B=S_A^{-1}B\hookrightarrow Q(B).
\end{align*}
特别地，在平坦情形下有亚纯函数的拉回 $\mathcal M(X)\longrightarrow\mathcal M(Y)$. 如果进一步假设 $f$ 忠实平坦，则此拉回是单射.
\end{numpara}
\begin{definition}\label{item:ag-1-1-4}
设 $X$ 是概形. $X$ 上的一个\newterm{ag-cartier-divisor}{Cartier 除子}{Cartier divisor}（本章简称除子）是层 $\mathcal M_X^\times/\mathcal O_X^\times$ 的一个整体截面. 如果它是子层 $(\mathcal O_X\cap\mathcal M_X^\times)/\mathcal O_X^\times$ 的整体截面，则称其为\newterm{ag-effective-divisor}{有效除子}{effective divisor}. 记
\begin{align*}
\operatorname{Div}(X)&:=\Gamma(X,\mathcal M_X^\times/\mathcal O_X^\times),\\
\operatorname{Div}^+(X)&:=\Gamma(X,(\mathcal O_X\cap\mathcal M_X^\times)/\mathcal O_X^\times),
\end{align*}
分别为 $X$ 上的除子全体及其中的有效除子全体.
\end{definition}
［具体地，一个除子 $D$ 由一个覆盖 $X=\bigcup_iU_i$ 和亚纯函数 $f_i\in\mathcal M(U_i)^\times$ 给出，并要求对任意 $i,j$，有 $(f_if_j)^{-1}\in\Gamma(U_i\cap U_j,\mathcal O_X)^\times$.］
群 $\operatorname{Div}(X)$ 是 Abel 群，按定义带有群同态
\begin{align*}
\mathcal M(X)^\times\xrightarrow{\operatorname{div}}\operatorname{Div}(X).
\end{align*}
不过，传统上将 $\operatorname{Div}(X)$ 写成加法群. 例如 $\operatorname{div}(fg)=\operatorname{div}(f)+\operatorname{div}(g)$. 此外，
\begin{align*}
\forall f,g\in\mathcal M(X),\qquad
\operatorname{div}(f)=\operatorname{div}(g)\ \Longleftrightarrow\ fg^{-1}\in\Gamma(X,\mathcal O_X^\times).
\end{align*}
子集 $\operatorname{Div}^+(X)$ 是 $\operatorname{Div}(X)$ 的子幺半群，满足 $\operatorname{Div}^+(X)\cap(-\operatorname{Div}^+(X))=\{0\}$. 因而由
\begin{align*}
D\leqslant D'\ \Longleftrightarrow\ (D'-D)\in\operatorname{Div}^+(X)
\end{align*}
定义的关系 $\leqslant$ 使 $\operatorname{Div}(X)$ 成为一个（偏）序群.
\begin{numpara}[可逆理想的回顾]\label{item:ag-1-1-5}
设 $X$ 是概形，$\mathcal L$ 是 $\mathcal O_X$-模. 如果每一点都有开邻域 $U$，使得 $\mathcal L|_U\simeq\mathcal O_U$，则称 $\mathcal L$ 局部自由且秩为 $1$.
此外，定义对偶 $\mathcal O_X$-模为层
\begin{align*}
\mathcal L^\vee:=\mathcal Hom_{\mathcal O_X}(\mathcal L,\mathcal O_X):
U\longmapsto\operatorname{Hom}_{\mathcal O_U}(\mathcal L|_U,\mathcal O_U).
\end{align*}
于是有典范的 $\mathcal O_X$-模态射
\begin{align*}
\mathcal L\otimes_{\mathcal O_X}\mathcal L^\vee\longrightarrow\mathcal O_X.
\end{align*}
如果 $\mathcal L$ 拟凝聚且此态射是同构，则称 $\mathcal L$ \newterm{ag-invertible-sheaf}{可逆}{invertible sheaf}.
可以证明，$\mathcal L$ 可逆当且仅当它局部自由且秩为 $1$.
\end{numpara}
\begin{definition}\label{item:ag-1-1-6}
$\mathcal M_X$ 的一个\newterm{ag-invertible-fractional-ideal}{可逆分式理想}{invertible fractional ideal}是 $\mathcal M_X$ 的一个 $\mathcal O_X$-子模，它作为 $\mathcal O_X$-模是可逆的.
\end{definition}
\emph{注：}回忆，我们已将``拟凝聚''条件包含在可逆性的定义中. 如果 $X=\operatorname{Spec}(A)$，则 $\mathcal M_X$ 的一个可逆分式理想可等同于 $Q(A)$ 的一个分式理想，后者作为 $A$-模可逆（也称``秩为 $1$ 的投射模''）. 回忆，这样的分式理想未必是主理想，只是``局部为主理想''.
\begin{proposition}\label{item:ag-1-1-7}
$\mathcal M_X$ 的 $\mathcal O_X$-子模 $\mathcal J$ 可逆，当且仅当每一点都有开邻域 $U$，使得 $\mathcal J|_U=\mathcal O_Uf$，其中 $f\in\mathcal M(U)^\times=\Gamma(U,\mathcal M_X^\times)$ 是亚纯函数.
\end{proposition}
\begin{proof}
取 $U$，使得存在同构 $\mathcal O_U\xrightarrow\sim\mathcal J|_U$. 单位截面 $1$ 在 $\mathcal M(U)^\times$ 中的像 $f$ 即符合要求.
\end{proof}
\begin{corollary}\label{item:ag-1-1-8}
沿用上一命题的记号，
\begin{enumerate}[label=\roman*)]
\item 可逆分式理想的逆 $\mathcal J^\vee$ 可典范地等同于唯一的分式理想 $\mathcal J^{-1}$，其定义条件为：对每个开集 $U$，若 $\mathcal J|_U=f\mathcal O_U$，则 $\mathcal J^{-1}=f^{-1}\mathcal O_U$.
\item 两个可逆分式理想的乘积 $\mathcal J_1\mathcal J_2$ 仍是可逆分式理想，且典范映射 $\mathcal J_1\otimes_{\mathcal O_X}\mathcal J_2\longrightarrow\mathcal J_1\mathcal J_2$ 是同构.
\end{enumerate}
\end{corollary}
特别地，$\mathcal M_X$ 的可逆分式理想全体 $\operatorname{Id.Fr.Inv}(X)$ 是一个 Abel 群. 同样，$\mathcal O_X$ 的可逆理想全体 $\operatorname{Id.Inv}(X)$ 是前者的一个交换子幺半群.
\begin{remark}\label{item:ag-1-1-9}
设 $\mathcal L$ 是可逆层. 截面 $s\in\Gamma(X,\mathcal L)$ 可等同于一个 $\mathcal O_X$-模态射 $\mathcal O_X\xrightarrow{\varphi_s}\mathcal L$，从而等同于态射 $\mathcal L^\vee\xrightarrow{\varphi_s^\vee}\mathcal O_X$. 如果 $\varphi_s$ 是单射，则称截面 $s$ 正则. 此时 $\varphi_s^\vee$ 给出 $\mathcal L^\vee$ 到 $\mathcal O_X$ 的一个（可逆）理想的同构. 容易验证，这一过程诱导双射
\begin{align*}
\left\{\begin{aligned}&\text{二元组 $(\mathcal L,s)$ 的同构类，}\\&\text{其中 $\mathcal L$ 是可逆层，}\\&\text{$s$ 是 $\mathcal L$ 的正则截面}\end{aligned}\right\}
\xrightarrow\sim\operatorname{Id.Inv}(X),
\end{align*}
这里，$(\mathcal L,s)$ 与 $(\mathcal L',s')$ 之间的同构是满足 $\varphi_{s'}\circ\alpha=\varphi_s$ 的同构 $\alpha:\mathcal L\xrightarrow\sim\mathcal L'$. 右边的运算当然来自张量积 $(\mathcal L\otimes_{\mathcal O_X}\mathcal L',s\otimes s')$.
将 $\mathcal L$ 的一个亚纯截面定义为层 $\mathcal L\otimes\mathcal M_X$ 的一个截面；如果由它得到的态射 $\mathcal O_X\longrightarrow\mathcal L\otimes\mathcal M_X$ 是单射，则称此截面正则. 与上面相同，有自然双射
\begin{align*}
\left\{\begin{aligned}&\text{二元组 $(\mathcal L,s)$ 的同构类，}\\&\text{其中 $\mathcal L$ 是可逆层，}\\&\text{$s$ 是 $\mathcal L$ 的正则亚纯截面}\end{aligned}\right\}
\xrightarrow\sim\operatorname{Id.Fr.Inv}(X).\tag*{$\diamond$}
\end{align*}
\renewcommand{\envenddiamond}{}% final marker is on the display
\end{remark}
\begin{numpara}\label{item:ag-1-1-10}
按如下方式定义群同态
\begin{align*}
\operatorname{Div}(X)&\longrightarrow\operatorname{Id.Fr.Inv}(X),\\
D&\longmapsto\mathcal I_X(D).
\end{align*}
设 $D\in\operatorname{Div}(X)$ 由覆盖 $X=\bigcup_iU_i$ 和亚纯函数 $f_i\in\mathcal M(U_i)^\times$ 定义. 对任意 $i,j$，条件 $(f_if_j)^{-1}\in\Gamma(U_i\cap U_j,\mathcal O_X)^\times$ 表明，$\mathcal M_{U_i\cap U_j}$ 的子层满足 $f_i\mathcal O_{U_i\cap U_j}=f_j\mathcal O_{U_i\cap U_j}$. 因而，$\mathcal M_{U_i}$ 的子层 $f_i\mathcal O_{U_i}$ 可粘合成 $\mathcal M_X$ 的子层 $\mathcal I_X(D)$；按构造，这是可逆分式理想.
\end{numpara}
\begin{proposition}\label{item:ag-1-1-11}
上述态射是 Abel 群的同构
\begin{align*}
\operatorname{Div}(X)\xrightarrow\sim\operatorname{Id.Fr.Inv}(X).
\end{align*}
它将有效除子映到 $\mathcal O_X$ 的可逆理想上，因而限制为幺半群同构
\begin{align*}
\operatorname{Div}^+(X)\xrightarrow\sim\operatorname{Id.Inv}(X).
\end{align*}
\end{proposition}
\begin{proof}
只需构造逆向双射. 给定可逆分式理想 $\mathcal J\subset\mathcal M_X$，选取开覆盖 $X=\bigcup_iU_i$ 和亚纯函数 $f_i\in\mathcal M(U_i)$，使得 $\mathcal J|_{U_i}=\mathcal O_{U_i}f_i$. 则对任意 $i,j$，等式 $\mathcal O_{U_i\cap U_j}f_i=\mathcal O_{U_i\cap U_j}f_j$ 表明 $f_if_j^{-1}\in\Gamma(U_i\cap U_j,\mathcal O_X^\times)$. 因而，该覆盖和这些 $f_i$ 给出 $\mathcal M_X^\times/\mathcal O_X^\times$ 的一个截面.
\end{proof}
用注~\ref{item:ag-1-1-9}的语言说，除子 $D$ 被映到二元组 $(\mathcal O_X(D),s_D)$，其中令 $\mathcal O_X(D):=\mathcal I_X(D)^\vee$，而 $s_D$ 是由 $\mathcal M_X$-模态射 $\mathcal M_X\longrightarrow\mathcal O_X(D)\otimes_{\mathcal O_X}\mathcal M_X$ 定义的亚纯截面. 此态射对偶于由包含 $\mathcal I_X(D)\subset\mathcal M_X$ 诱导的态射 $\mathcal I_X(D)\otimes_{\mathcal O_X}\mathcal M_X\longrightarrow\mathcal M_X$.
于是有典范同构 $\mathcal O_X(D+D')\xrightarrow\sim\mathcal O_X(D)\otimes\mathcal O_X(D')$，它将 $s_{D+D'}$ 映到 $s_D\otimes s_{D'}$.
\begin{numpara}[有效除子所对应的子概形]\label{item:ag-1-1-12}
设 $D$ 是 $X$ 上的有效除子. 对应的分式理想 $\mathcal I_X(D)$ 是 $\mathcal O_X$ 的一个（通常意义下的）理想，因而定义 $X$ 的一个闭子概形 $Y(D)$. 局部地，这个子概形由一个方程定义，且此方程正则. 称 $Y(D)$ 到 $X$ 的浸入是余维 $1$ 的正则浸入（后面将介绍一般的正则浸入的定义）. 因而有双射
\begin{align*}
\operatorname{Div}^+(X)\xrightarrow\sim
\{\text{$X$ 的余维 $1$ 的正则浸入闭子概形}\}.
\end{align*}
有时也滥用术语，将这样的子概形称为``除子''.
\end{numpara}
\section{除子的线性等价与可逆层}\label{sec:ag-1-2}
\begin{numpara}[Picard 群]\label{item:ag-1-2-1}
设 $X$ 是概形. 记 $\operatorname{Pic}(X)$ 为可逆层的同构类全体. 张量积使 $\operatorname{Pic}(X)$ 成为 Abel 群，称为\newterm{ag-picard-group}{Picard 群}{Picard group}. 有群同态
\begin{align*}
l:\operatorname{Div}(X)\xrightarrow\sim\operatorname{Id.Fr.Inv}(X)\longrightarrow\operatorname{Pic}(X),
\end{align*}
它将可逆分式理想映到其逆作为可逆层的同构类. 用注~\ref{item:ag-1-1-9}的语言说，此同态就是忘掉截面 $s$.
\end{numpara}
\begin{numpara}[上同调解释]\label{item:ag-1-2-2}
有典范同构
\begin{align*}
\operatorname{Pic}(X)\xrightarrow\sim\check H^1(X,\mathcal O_X^\times)
\qquad\text{（\v Cech 上同调）}.
\end{align*}
熟悉这一语言之后，这实质上是不言而喻的. 具体可作如下描述. 固定可逆层 $\mathcal L$，选取覆盖 $X=\bigcup_iU_i$ 以及同构 $\varphi_i:\mathcal L|_{U_i}\xrightarrow\sim\mathcal O_{U_i}$. 于是对任意 $i,j$，有函数 $f_{ij}=\varphi_j\circ\varphi_i^{-1}\in\Gamma(U_i\cap U_j,\mathcal O_X^\times)$，并对任意 $i,j,k$ 满足上闭链条件 $f_{jk}f_{ij}=f_{ik}$. 因而它定义一个 $1$-上闭链，进而定义 $H^1(X,\mathcal O_X^\times)$ 中的上同调类. 反过来，给定这样的类，用一个 \v Cech $1$-上闭链 $((U_i)_{i\in I},(f_{ij})_{i,j\in I})$ 表示它，其中 $f_{ij}\in\Gamma(U_i\cap U_j,\mathcal O_X^\times)$. 于是 $f_{ij}$ 定义同构 $(\mathcal O_{U_i})|_{U_i\cap U_j}\xrightarrow\sim(\mathcal O_{U_j})|_{U_i\cap U_j}$，而上闭链条件使这些层可粘合为一个 $\mathcal O_X$-模，且它必然可逆.
回忆，对每个 $i$ 及每个 Abel 层 $\mathcal F$，都有典范映射 $\check H^i(X,\mathcal F)\longrightarrow H^i(X,\mathcal F)=R^i\Gamma(X,\mathcal F)$（$\Gamma(X,-)$ 的第 $i$ 个导出函子）；当 $i=1$ 时，此映射是同构（参见 Hartshorne，习题 III.4.4）. 考虑层的正合列
\begin{align*}
1\longrightarrow\mathcal O_X^\times\longrightarrow\mathcal M_X^\times
\longrightarrow\mathcal M_X^\times/\mathcal O_X^\times\longrightarrow1.
\end{align*}
它给出上同调长正合列
\begin{align*}
\Gamma(X,\mathcal O_X^\times)&\hookrightarrow\Gamma(X,\mathcal M_X^\times)
\longrightarrow\Gamma(X,\mathcal M_X^\times/\mathcal O_X^\times)\\
&\xrightarrow\delta H^1(X,\mathcal O_X^\times)\longrightarrow H^1(X,\mathcal M_X^\times).
\end{align*}
用已经引入的记号，可将它改写为
\begin{align*}
\Gamma(X,\mathcal O_X^\times)\hookrightarrow\mathcal M(X)^\times
\longrightarrow\operatorname{Div}(X)\xrightarrow\delta\operatorname{Pic}(X)
\longrightarrow H^1(X,\mathcal M_X^\times).
\end{align*}
读者可以练习验证，上述态射 $\delta$ 就是上一分目中记为 $l$ 的态射.
\end{numpara}
\begin{definition}\label{item:ag-1-2-3}
如果 $X$ 上的除子 $D$ 具有形式 $\operatorname{div}(f)$，其中 $f\in\mathcal M(X)$ 是亚纯函数，则称 $D$ 是\newterm{ag-principal-divisor}{主除子}{principal divisor}. 如果两个除子 $D_1,D_2$ 之差 $D_1-D_2$ 是主除子，则称它们\newterm{ag-linearly-equivalent}{线性等价}{linearly equivalent}.
\end{definition}
\emph{注：}除子 $D$ 是主除子，当且仅当对应的可逆分式理想 $\mathcal I_X(D)\subset\mathcal M_X$ 具有形式 $f\mathcal O_X$，其中 $f$ 是 $\mathcal M_X$ 的一个整体正则截面. 因此，这里的``主''与环的理想的通常术语相容.
由前面的讨论，映射 $\operatorname{div}$ 诱导同构
\begin{align*}
\mathcal M(X)^\times/\Gamma(X,\mathcal O_X^\times)\xrightarrow\sim\operatorname{Div.Pr}(X),
\end{align*}
而映射 $l$ 诱导单态射
\begin{align*}
\operatorname{Div}(X)/\operatorname{Div.Pr}(X)\hookrightarrow\operatorname{Pic}(X).
\end{align*}
注意，当 $X$ 是整概形时，常值层 $\mathcal M_X$ 是松弛层. 因而 $H^1(X,\mathcal M_X^\times)=0$，所以上述态射是双射. 不假设整性时，有如下结果.
\begin{theorem}\label{item:ag-1-2-4}
设 $X$ 是 Noether 环上的拟射影概形. 则态射 $l$ 是满射，因而诱导同构
\begin{align*}
\operatorname{Div}(X)/\operatorname{Div.Pr}(X)\xrightarrow\sim\operatorname{Pic}(X).
\end{align*}
\end{theorem}
\begin{proof}
参见 EGA IV (21.3.5). 问题在于：给定一个可逆层 $\mathcal L$，为它找一个正则亚纯截面，即 $\mathcal L\otimes_{\mathcal O_X}\mathcal M_X$ 的正则截面. 在仿射情形 $X=\operatorname{Spec}(A)$ 下，$\mathcal L$ 对应于可逆 $A$-模 $L$，而 $\mathcal L\otimes_{\mathcal O_X}\mathcal M_X$ 对应于可逆 $Q(A)$-模 $L\otimes_AQ(A)$. 但 $Q(A)$ 是 Artin 环，在这样的环上，每个秩为 $1$ 的投射模都是自由模. 因此，任意同构 $Q(A)\xrightarrow\sim L\otimes_AQ(A)$ 都给出所需的截面.
一般情形可归结到仿射情形. 借此机会，我们引入概形的\newterm{ag-schematically-dense}{概形论稠}{schematically dense}开集的概念. 设 $U$ 是概形 $X$ 的开集. 如果下列两个等价性质之一成立，就称 $U$ 在 $X$ 中概形论稠：
\begin{enumerate}[label=\roman*)]
\item 对 $X$ 的每个开集 $V$，限制映射 $\Gamma(V,\mathcal O_V)\longrightarrow\Gamma(U\cap V,\mathcal O_V)$ 是单射.
\item 对每个分离概形 $Y$ 和每个开集 $V$，限制映射 $\operatorname{Hom}(V,Y)\longrightarrow\operatorname{Hom}(V\cap U,Y)$ 是单射.
\end{enumerate}
容易验证，这样的开集必然在拓扑意义下稠密. 当 $X$ 既约时，反过来也成立. 事实上，若 $X=\operatorname{Spec}(A)$，则主开集 $D_f$ 概形论稠当且仅当关于 $f$ 的局部化映射是单射，也即当且仅当 $f$ 正则. 这引导我们考虑 $A$ 的\newterm{ag-associated-prime-ideal}{结合素理想}{associated prime ideal}，即作为 $A$ 的某个元素的零化子的素理想. 如果 $A$ 是 Noether 环，这些理想组成的集合 $\operatorname{Ass}(A)$ 有限. 此外，结合素理想的并恰好是正则元素集合的补集. 于是容易证明，$\operatorname{Spec}(A)$ 的开集 $U$ 概形论稠，当且仅当它包含 $\operatorname{Ass}(A)$. 当 $A$ 既约时，结合素理想对应于不可约分支的泛点，因而稠密与概形论稠等价.
下面给出一个稠密但不概形论稠的开集的例子. 设 $A$ 是 Noether 环，$x$ 是 $\operatorname{Spec}(A)$ 的闭点，对应理想为 $\mathfrak m$. 构造环 $A':=A\oplus\mathfrak m/\mathfrak m^2$，乘法为 $(a,m)(a',m')=(aa',am'+a'm)$. 显然的态射 $A'\longrightarrow A$ 诱导 $\operatorname{Spec}(A)$ 到 $\operatorname{Spec}(A')$ 的同胚. 现在，$\mathfrak m/\mathfrak m^2$ 中的每个元素 $m$ 都给出 $\Gamma(\operatorname{Spec}(A'),\mathcal O_{\operatorname{Spec}(A')})=A$ 的一个非零截面，它在开集 $U=\operatorname{Spec}(A')\setminus\{x\}$ 上为零，因而此开集不概形论稠. 例如取 $A=k[T]$，其中 $k$ 是域，$x$ 为 $T=0$ 定义的点. 此时环 $A'$ 同构于 $k[T,U]/(U^2,TU)$. 称 $x$ 是 $\operatorname{Spec}(A')$ 的一个\newterm{ag-embedded-point}{嵌入点}{embedded point}. 因此 $\operatorname{Spec}(A')$ 是一条``在原点加厚''的仿射直线.
在更一般的概形 $X$ 上，除子的补集总是概形论稠的. 如果 $X$ 是 Noether 概形，记 $\operatorname{Ass}(\mathcal O_X)$ 为满足如下条件的点 $x$ 所成的有限集合：$\mathcal O_{X,x}$ 的极大理想 $\mathfrak m_x$ 只含零因子. 归结到仿射情形，可证明（参见 EGA IV.3.1.8），开集 $U$ 在 $X$ 中概形论稠当且仅当它包含 $\operatorname{Ass}(\mathcal O_X)$.
回到定理. 当 $X$ 满足定理的假设时，已知每个有限集都包含在某个仿射开集中（这样的开集可取为某个超平面截面的补集）. 将此应用于 $\operatorname{Ass}(\mathcal O_X)$，便得到一个概形论稠的仿射开集 $U$. 根据已经处理的仿射情形，存在正则截面 $s_U\in\Gamma(U,\mathcal L\otimes_{\mathcal O_X}\mathcal M_X)$. 设 $V$ 是满足 $\mathcal L|_V\simeq\mathcal O_V$ 的仿射开集，并选定这样的同构. 截面 $(s_U)|_{U\cap V}$ 对应于亚纯函数 $f_{UV}\in\Gamma(U\cap V,\mathcal M_X^\times)$. 因为 $U\cap V$ 在 $V$ 中概形论稠，限制映射 $\mathcal M(V)\longrightarrow\mathcal M(U\cap V)$ 是双射. 所以 $f_{UV}$ 可唯一延拓为 $f_V\in\mathcal M(V)^\times$，后者对应于 $\mathcal L|_V$ 的正则亚纯截面 $s_V$. 用上述开集 $V$ 覆盖 $X$，就可将这些 $s_V$ 粘合成 $\mathcal L$ 的一个正则亚纯截面.
\end{proof}
\begin{numpara}[函子性]\label{item:ag-1-2-5}
对态射 $Y\xrightarrow f X$，可逆层的逆像 $f^*\mathcal L$ 仍是可逆层，因而有态射 $\operatorname{Pic}(X)\longrightarrow\operatorname{Pic}(Y)$. 但是，除子的逆像并不总有定义（考虑 $f$ 是一个点的闭浸入的情形）.
\end{numpara}
\begin{proposition}\label{item:ag-1-2-6}
在下列两种情形下，有``逆像''同态 $\operatorname{Div}(X)\longrightarrow\operatorname{Div}(Y)$：
\begin{enumerate}[label=\roman*)]
\item $f$ 平坦.
\item $X,Y$ 是整概形，且 $f(X)$ 包含一个开集. 特例包括：$f$ 双有理（即诱导 $X$ 的一个开集到 $Y$ 的一个开集的同构），以及 $f$ 是域上代数簇之间的支配态射.
\end{enumerate}
\end{proposition}
\Needspace{5\baselineskip}
\begin{proof}
\begin{enumerate}[label=\roman*)]
\item 设 $D$ 是除子，将它看作分式理想 $\mathcal J\hookrightarrow\mathcal M_X$. 因为 $f$ 平坦，有 $f^*\mathcal J\hookrightarrow f^*\mathcal M_X$，并且有单态射 $f^*\mathcal M_X\hookrightarrow\mathcal M_Y$，参见分目~\ref{item:ag-1-1-3}中的注. 所以 $f^*\mathcal J$ 是 $\mathcal M_Y$ 的分式理想，并定义 $Y$ 上的除子 $f^*(D)$.
\item 由假设，$f$ 的像包含 $X$ 的一个稠密开集 $U$. 设 $\varphi\in\mathcal M(X)$. 存在开集 $U_\varphi\subset U$，使得 $\varphi|_{U_\varphi}\in\mathcal O_X(U_\varphi)$. 因而
\begin{align*}
f^*(\varphi|_{U_\varphi})\in\mathcal O_Y(f^{-1}(U_\varphi))
\subset\mathcal M(U_\varphi)=\mathcal M(Y).
\end{align*}
这样便定义了拉回 $\mathcal M(X)\xrightarrow{f^*}\mathcal M(Y)$. 再设除子 $D$ 由覆盖 $X=\bigcup_iU_i$ 和 $\varphi_i\in\mathcal M(X)^*$ 给出. 逆像覆盖 $Y=\bigcup_if^{-1}(U_i)$ 及函数 $f^*(\varphi_i)$ 定义 $Y$ 上的除子 $f^*D$.\qedhere
\end{enumerate}
\end{proof}
\section{余维 \texorpdfstring{$1$}{1} 的圈（Weil 除子）}\label{sec:ag-1-3}
\begin{numpara}[圈]\label{item:ag-1-3-1}
设 $X$ 是 Noether 概形. $X$ 中的一个\newterm{ag-cycle}{圈}{cycle}是 $X$ 的底空间的一个不可约闭子集. 因而圈与 $X$ 的点一一对应（对一个圈取其泛点，对一个点 $x$ 取其闭包 $\overline{\{x\}}$）. 如果 $\dim(\mathcal O_{X,x})=p$，则称圈 $\overline{\{x\}}$ 的余维为 $P$.
记 $\mathfrak Z(X)$、$\mathfrak Z^1(X)$ 分别为以 $X$ 的圈、余维 $1$ 的圈为基的自由 Abel 群.
如果 $X=\operatorname{Spec}(A)$，则余维 $1$ 的圈对应于高度 $1$ 的素理想.
\end{numpara}
\begin{numpara}\label{item:ag-1-3-2}
我们希望定义一个同态 $\operatorname{Div}(X)\longrightarrow\mathfrak Z^1(X)$. 先定义
\begin{align*}
\operatorname{Div}^+(X)&\longrightarrow\mathfrak Z^1(X),\\
D&\longmapsto\operatorname{cyc}(D):=
\sum_{x\in X^1}\operatorname{long}_{\mathcal O_{X,x}}(\mathcal O_{Y(D),x})[x].
\end{align*}
注意，因为 $Y(D)$ 的余维为 $1$，所以点 $x\in X^1\cap Y(D)$ 是 $Y(D)$（的底空间）的某个不可约分支的泛点. 因而上面的和确实是有限和. 此外，对这样的泛点，局部环 $\mathcal O_{Y(D),x}$ 是 Artin 环，且上式中的长度就是它作为 Artin 环的长度.
验证 $\operatorname{cyc}(D+D')=\operatorname{cyc}(D)+\operatorname{cyc}(D')$. 设 $f,f'$ 分别是 $Y(D),Y(D')$ 在 $x$ 附近的局部方程，所要证明的是
\begin{align*}
\operatorname{long}(\mathcal O_{X,x}/(f))+
\operatorname{long}(\mathcal O_{X,x}/(f'))
=\operatorname{long}(\mathcal O_{X,x}/(ff')).
\end{align*}
这由下列正合列得出：
\begin{align*}
\mathcal O_{X,x}/(f)\overset{\times f'}{\hookrightarrow}
\mathcal O_{X,x}/(ff')\twoheadrightarrow\mathcal O_{X,x}/(f').
\end{align*}
该列正合，因为 $f'$ 是正则元素.
\end{numpara}
\begin{proposition}\label{item:ag-1-3-3}
上述同态可典范地延拓为同态 $\operatorname{Div}(X)\longrightarrow\mathfrak Z^1(X)$.
\end{proposition}
\begin{proof}
若已知 $\operatorname{Div}^+(X)$ 生成群 $\operatorname{Div}(X)$，则结论显然. 当 $X$ 是整仿射概形时，情况正是如此. 事实上，除子 $D$ 由有限覆盖 $X=\bigcup_iU_i$ 和亚纯函数 $f_i\in\mathcal M(U_i)=\mathcal M(X)$ 给出. 于是可找到函数 $f$，使得 $ff_i\in\Gamma(U_i,X)$，从而 $D=D'-\operatorname{div}(f)$ 是两个有效除子之差.
一般情形下，只需对 $x\in X^1$ 延拓映射 $D\mapsto\operatorname{mult}_x(D):=\operatorname{long}(\mathcal O_{Y(D),x})$. 可以找到 $x$ 的仿射开邻域 $U$ 和亚纯函数 $f_U\in\mathcal M(U)$，使得 $D|_U=\operatorname{div}(f_U)$. 按定义，$f_U=f_{U,+}f_{U,-}^{-1}$，其中 $f_{U,\pm}\in\Gamma(U,\mathcal O_X)$ 是正则函数. 于是令
\begin{align*}
\operatorname{mult}_x(D):=
\operatorname{mult}_x(\operatorname{div}(f_{U,+}))-
\operatorname{mult}_x(\operatorname{div}(f_{U,-})).\tag*{$\square$}
\end{align*}
\renewcommand{\qed}{}% final marker is on the display
\end{proof}
\begin{remark}\label{item:ag-1-3-4}
假设 $X$ 是整概形，$x$ 是余维 $1$ 的正则点. 则 $\mathcal O_{X,x}$ 是离散赋值环. 记 $v_x:\mathcal M(X)\longrightarrow\mathbb Z$ 为函数域 $\mathcal M(X)$ 上对应的赋值. 如果 $D$ 在 $x$ 附近局部具有形式 $\operatorname{div}(f)$，则 $\operatorname{mult}_x(D)=v_x(f)$.
\end{remark}
\begin{numpara}\label{item:ag-1-3-5}
记 $\mathfrak Z_{\mathrm{prin}}^1(X)$ 为 $\mathfrak Z^1$ 中由形如 $\operatorname{cyc}(\operatorname{div}(f))$、$f\in\mathcal M(X)$ 的元素组成的子群（称这些元素为主圈）. 因而有同态
\begin{align*}
\operatorname{Div}(X)/\operatorname{Div.Pr}(X)
\longrightarrow\mathfrak Z^1(X)/\mathfrak Z_{\mathrm{prin}}^1(X).
\end{align*}
右边的商称为``$X$ 的\newterm{ag-class-group}{类群}{class group}''，记为 $\mathcal Cl(X)$.
\emph{例：}如果 $A$ 是 Dedekind 环，例如某数域的整数环，则 $\mathcal Cl$ 就是算术中考虑的类群. 回忆，这个 $\mathcal Cl$ 控制着 $A$ 偏离主理想环的程度，以及分解为不可约元素时偏离唯一性的程度. 更一般地，有如下结果.
\end{numpara}
\begin{proposition}\label{item:ag-1-3-6}
一个（Noether）整环是唯一分解整环，当且仅当它整闭，且 $\mathcal Cl(\operatorname{Spec}(A))=0$.
\end{proposition}
回忆，整环称为\newterm{ag-unique-factorization-domain}{唯一分解整环}{unique factorization domain}，如果每个非零元素都有``唯一''的不可约元素分解. 已知这等价于所有高度 $1$ 的素理想都是主理想.
\begin{proof}
我们承认唯一分解整环必正规（这不难）. 假设 $A$ 是唯一分解整环. $\operatorname{Spec}(A)$ 的余维 $1$ 的圈 $Z$ 就是 $A$ 的一个高度 $1$ 的理想. 这样的理想是由某函数 $f$ 生成的主理想，所以 $\operatorname{div}(f)=1.Z$，从而 $\mathcal Cl(\operatorname{Spec}(A))=0$.
反过来，假设 $\mathcal Cl(\operatorname{Spec}(A))=0$. 设 $\mathfrak p$ 是高度 $1$ 的理想，$Z$ 是对应的圈. 我们要证明 $\mathfrak p$ 为主理想. 由假设，存在 $f\in Q(A)$ 使得 $\operatorname{div}(f)=1.Z$. 因为 $A$ 正规，故在余维 $1$ 处正则，所以对每个高度 $1$ 的 $\mathfrak p'$，$A_{\mathfrak p'}$ 都是离散赋值环. 因而，当 $\mathfrak p'\ne\mathfrak p$ 时 $v_{\mathfrak p'}(f)=0$，而 $v_{\mathfrak p}(f)=1$. 特别地，对所有高度 $1$ 的 $\mathfrak p'$，都有 $f\in A_{\mathfrak p'}$. 已知正规（即整闭）环具有如下性质：
\begin{align*}
A=\bigcap_{\operatorname{ht}(\mathfrak p)=1}A_{\mathfrak p}.
\end{align*}
由此得到 $f\in A\cap\mathfrak pA_{\mathfrak p}$. 证明 $f$ 生成 $\mathfrak p$. 设 $g\in\mathfrak p$，再次应用上述结果，便得到 $g/f\in A$.
\end{proof}
\begin{theorem}\label{item:ag-1-3-7}
假设 $X$ 是正规整概形. 则同态 $\operatorname{cyc}:\operatorname{Div}(X)\longrightarrow\mathfrak Z^1(X)$ 是单射. 此外，它是双射当且仅当 $X$ 的所有局部环都是唯一分解整环.
\end{theorem}
\begin{proof}
回忆，$X$ 正规蕴含余维 $1$ 的点处的局部环正则. 为证明单射性，只需证明
\begin{align*}
\operatorname{cyc}^{-1}(\mathfrak Z^{1,+}(X))=\operatorname{Div}^+(X).
\end{align*}
设 $D$ 是满足对每个 $x$ 都有 $\operatorname{mult}_x(D)\geqslant0$ 的除子. 取仿射开集 $U=\operatorname{Spec}(A)$，使得其上 $D=\operatorname{div}(f)$. 对 $x\in U$，有 $\operatorname{mult}_x(D)=v_x(f)\geqslant0$，因而 $f\in\mathcal O_{X,x}$. 因为 $A$ 正规，得到 $f\in A$. 对上述每个 $U$ 都如此，故 $D\geqslant0$.
再证满射性. 先假设 $X$ 的局部环都是唯一分解整环. 设 $Z\in\mathfrak Z^1(X)$. 对点 $x\in X$，记 $Z_x\in\mathfrak Z^1(\operatorname{Spec}(\mathcal O_{X,x}))$ 为 $Z$ 在 $\operatorname{Spec}(\mathcal O_{X,x})$ 上的限制. 因为 $\mathcal O_{X,x}$ 是唯一分解整环，存在 $f_x\in\mathcal M_{X,x}=\mathcal M(X)$，使得 $D_x=\operatorname{cyc}(\operatorname{div}(f_x))$. 因而 $\mathfrak Z^1(X)$ 中的 $D-\operatorname{cyc}(\operatorname{div}(f_x))$ 的支集不与 $x$ 相交，所以存在开集 $U$，其上仍有 $D|_U=\operatorname{cyc}(\operatorname{div}(f_x)|_U)$. 用这样的开集 $U$ 覆盖 $X$，便得到一个除子.
反过来，假设同态 $\operatorname{cyc}$ 满射. 设 $x\in X$，且 $Z$ 是经过 $x$ 的圈. 写成 $Z=\operatorname{cyc}(D)$. 在 $x$ 附近局部地，有 $D=\operatorname{div}(f)$，其中 $f$ 是某个函数. 因而 $Z$ 在 $\mathcal O_{X,x}$ 中是主的.
\end{proof}
\begin{corollary}\label{item:ag-1-3-8}
如果 $X$ 局部为唯一分解整环，且在某 Noether 环上拟射影，则有同构 $\operatorname{Pic}(X)\xrightarrow\sim\mathcal Cl(X)$.
\end{corollary}
\begin{remark}\label{item:ag-1-3-9}
在上述推论的假设下，对 $X$ 的任意开集 $U$，限制同态 $\operatorname{Pic}(X)\longrightarrow\operatorname{Pic}(U)$ 都是满射.
\end{remark}
\begin{example}\label{item:ag-1-3-10}
如果 $X$ 是代数闭域上的非奇异射影曲线，则余维 $1$ 的圈就是闭点. 对每个有限态射 $f:Y\longrightarrow X$，其中 $Y$ 为非奇异曲线，可如下自然地诱导同态 $\mathfrak Z^1(X)\longrightarrow\mathfrak Z^1(Y)$：若 $t$ 是 $\mathcal O_{X,x}$ 中 $\mathfrak m_x$ 的生成元，则令 $f^*(x)=\sum_{f(y)=x}v_y(f^*t)$，其中 $f^*t\in\mathcal O_{Y,y}$. 这与 $t$ 的选择无关，然后按线性延拓. 可以验证，它保持线性等价（注意，因为 $f$ 支配，拉回 $\mathcal M(X)\longrightarrow\mathcal M(Y)$ 有良好定义，并使 $\mathcal M(Y)$ 成为 $\mathcal M(X)$ 的有限扩张），因而得到同态 $\mathcal Cl(X)\xrightarrow{f^*}\mathcal Cl(Y)$. 读者可以练习验证，通过推论~\ref{item:ag-1-3-8}的同构，此同态对应于拉回 $\operatorname{Pic}(X)\longrightarrow\operatorname{Pic}(Y)$.
\end{example}
\begin{theorem}[Auslander--Buschbaum]\label{item:ag-1-3-11}
Noether 正则局部环是唯一分解整环.
\end{theorem}
\begin{proof}
参见 EGA IV.21.11.
\end{proof}
\begin{proposition}\label{item:ag-1-3-12}
设 $X$ 是正规整概形，$Z\subset X$ 是真闭子集，补集为 $U$. 满的限制同态 $\mathcal Cl(X)\longrightarrow\mathcal Cl(U)$ 在 $Z$ 的余维 $\geqslant2$ 时是同构；在 $Z$ 是余维 $1$ 的圈时，诱导正合列
\begin{align*}
\mathbb Z\longrightarrow\mathcal Cl(X)\longrightarrow\mathcal Cl(U)\longrightarrow0,
\end{align*}
其中左边的映射定义为 $n\mapsto n.Z$.
\end{proposition}
\begin{proof}
显然.
\end{proof}
\begin{numpara}[应用]\label{item:ag-1-3-13}
$\mathcal Cl(\mathbb P_k^n)=\mathbb Z$. 事实上，写成 $\mathbb P_k^n=\mathbb A_k^n\sqcup\mathbb P_k^{n-1}$，并利用 $K[X_1,\ldots,X_n]$ 是唯一分解整环. （$\mathbb P^n$ 中超曲面的）次数映射诱导同态 $\mathcal Cl(\mathbb P_k^n)\longrightarrow\mathbb Z$，它是 $n\mapsto n[\mathbb P_k^{n-1}]$ 的逆映射.
\end{numpara}
\begin{proposition}\label{item:ag-1-3-14}
设 $X$ 是正规整概形. 则有同构 $\mathcal Cl(X)\xrightarrow\sim\mathcal Cl(X\times\mathbb A^1)$.
\end{proposition}
\begin{numpara}[余维 $1$ 的圈所对应的子概形]\label{item:ag-1-3-15}
设 $Z=\sum n_xx$ 是 $X$ 上的除子. 为它对应一个子概形 $Y(Z)$. 局部地，如果 $X=\operatorname{Spec}(A)$，则定义 $Y(Z)$ 的理想是映射 $A\longrightarrow\prod_xA_{\mathfrak p_x}/\mathfrak p_x^{n_x}$ 的核.
\end{numpara}
